\documentclass[../main.tex]{subfiles}
\begin{document}
\section{Protección de Datos}

Esta unidad se mira desde la perspectiva del \textbf{CISO} (\emph{Chief Information
Security Officer}), el responsable de velar por la seguridad de la información de la
empresa (con implicancias legales si hay un incidente severo). Se aplican todos los
principios ya vistos para proteger la \textbf{CIA} (Confidencialidad, Integridad y
Disponibilidad) de los datos.

%======================================================================
\subsection{El problema y el ciclo de protección}
%======================================================================

El problema motivador: en una empresa se define una política de protección de datos
personales (ej.\ ``un empleado no puede ver los sueldos de otro''). Los \textbf{controles
de acceso son parte de la solución pero no son suficientes}: un mismo dato (un sueldo)
fluye por RR.HH.\ $\to$ App $\to$ Excel $\to$ Email $\to$ nube, con backups en el medio.
Hay que protegerlo en todo su recorrido.

\begin{definicion}{Dato personal}
Todo lo que permite \textbf{identificar a una persona} solamente por ese mismo dato:
nombre y apellido, dirección de mail, número de documento, domicilio. Dentro de la
organización hay además datos \textbf{sensibles} (ej.\ el sueldo) que no se quieren
divulgar.
\end{definicion}

\textbf{Ciclo de protección de la información} (es un ciclo, vuelve a empezar):
\begin{enumerate}[nosep]
  \item \textbf{Gobierno (Govern):} gestión, toma de decisiones.
  \item \textbf{Descubrimiento (Discover):} qué datos hay en el universo de la organización.
  \item \textbf{Protección (Protect):} cómo y qué datos protegemos (no todos son sensibles).
  \item \textbf{Cumplimiento (Comply):} asegurar que la protección efectivamente ocurre.
  \item \textbf{Detección (Detect):} detectar filtraciones / \emph{breaches}.
  \item \textbf{Respuesta (Respond):} responder ante filtraciones y evitar que se repitan.
\end{enumerate}

\subsubsection{Data Governance}
Son \textbf{todas las tareas para asegurar que los datos estén seguros, privados,
íntegros y disponibles}. Incluye las medidas que toman las \textbf{personas}, los
\textbf{procesos} que siguen y la \textbf{tecnología} que los respalda durante el
\textbf{ciclo de vida de los datos}. Define los roles dentro de la organización.
Establece estándares internos (\textbf{políticas de datos}) sobre: recopilación,
almacenamiento, procesamiento, retención y eliminación.

\textbf{Data Life Cycle (PwC):} (1) Create \& collect $\to$ (2) Store \& transmit $\to$
(3) Use \& distribute $\to$ (4) Retain $\to$ (5) Dispose and destroy (al vencer el
período de retención).

\subsubsection{Data Roles (jerarquía)}
\begin{itemize}[nosep]
  \item \textbf{Data Trustee:} posición estratégica/\emph{cabinet-level}, responsable
        \textbf{legalmente} de los datos de su unidad.
  \item \textbf{Data Steward:} responsable interno que \textbf{aprueba accesos} a los
        datos dentro de su dominio.
  \item \textbf{Data Custodian:} autorizado por un Steward a \textbf{otorgar acceso}
        a los usuarios dentro del dominio del Steward.
  \item \textbf{Data User:} quien efectivamente \textbf{accede} y usa los datos.
\end{itemize}
A medida que se baja en la jerarquía, el dominio de datos manejado es cada vez más chico.

%======================================================================
\subsection{Marco legal y regulaciones}
%======================================================================

\begin{destacado}
La distinción más preguntable: la \textbf{Ley argentina 25.326} regula las bases de
datos (públicas y privadas). El \textbf{Derecho al Olvido} lo regula el \textbf{GDPR
europeo}, NO la ley argentina.
\end{destacado}

\subsubsection{Regulación argentina: Ley 25.326}
\begin{definicion}{Ley 25.326 -- Protección de Datos Personales}
Ley argentina (de los años 80--90). Aplica a todo sistema que opera en Argentina.
\begin{itemize}[nosep]
  \item \textbf{Art.\ 1 (Objeto):} protección integral de los datos personales asentados
        en archivos, registros, bancos de datos u otros medios técnicos de tratamiento,
        \textbf{sean públicos o privados}, para garantizar el derecho al honor y a la
        intimidad de las personas (conforme art.\ 43, párr.\ 3 de la Constitución).
  \item \textbf{Art.\ 9 (Seguridad de los datos):} el responsable/usuario del archivo
        debe adoptar las medidas técnicas y organizativas necesarias para garantizar
        seguridad y confidencialidad, evitando adulteración, pérdida, consulta o
        tratamiento no autorizado, y permitiendo detectar desviaciones. Queda
        \textbf{prohibido} registrar datos en bancos que no reúnan condiciones técnicas
        de integridad y seguridad.
\end{itemize}
\end{definicion}

\begin{destacado}
La ley \textbf{no especifica una tecnología} concreta. Esto es a propósito (``para
bien''): la generalidad hace que la ley \textbf{siga siendo aplicable hoy} pese a ser
de hace décadas. Implica también que las personas pueden \textbf{ver, modificar e
incluso pedir el borrado} de sus datos (hábeas data).
\end{destacado}

\paragraph{Hábeas data vs.\ hábeas corpus.}
El \textbf{hábeas data} es el derecho por el cual sos dueño de tus datos: podés pedir
verlos, rectificarlos o que se borren (a diferencia del \emph{hábeas corpus}, que
protege la libertad física de la persona). Empresas como Facebook/Twitter permiten
descargar todo tu historial y borrar la cuenta.

\subsubsection{Regulaciones internacionales}
\begin{center}
\begin{tabular}{@{}>{\bfseries}l p{10.5cm}@{}}
\toprule
Norma & Descripción \\
\midrule
HIPAA (1996) & \emph{Health Insurance Portability and Accountability Act}. Protección
  de datos de salud (PHI/PII). Surge la \textbf{anonimización} de datos (operar sobre
  historias clínicas sin saber de quién son). \\
GDPR (2018) & \emph{General Data Protection Regulation} (Unión Europea). La más conocida.
  Aplica aunque el sistema esté en Argentina si hay usuarios/infraestructura en la UE.
  Pide \textbf{consentimiento explícito} (de ahí los \emph{banners} de cookies). Regula
  el \textbf{Derecho al Olvido}. \\
PCI-DSS & \emph{Payment Card Industry Data Security Standard}. Tarjetas de crédito.
  \textbf{Acotó la cantidad de datos a proteger} (número, titular, vencimiento, código
  de seguridad/CVV, banda magnética). Introdujo la \textbf{tokenización}. \\
SOX (2002) & \emph{Sarbanes-Oxley Act} (ley federal de EE.UU.). Empresas que cotizan
  en bolsa; exige transparencia. Retención de auditoría \textbf{7 años}. \\
FFIEC & \emph{Assessment Tool} (framework, hoy en desuso). \\
FDA Cybersecurity & Ciberseguridad en dispositivos médicos. \\
\bottomrule
\end{tabular}
\end{center}

\begin{destacado}
El gran aporte de \textbf{PCI-DSS} (``Blue Print''): a la hora de proteger datos hay
que \textbf{acotar la cantidad de datos en base a su criticidad}. No proteger todo por
igual, sino concentrarse en lo crítico.
\end{destacado}

\paragraph{Tokenización (PCI-DSS).} Los datos críticos de la tarjeta se guardan ``bajo
7 llaves'' en una base recontra protegida; el sistema opera con un \textbf{token} que
apunta a la tarjeta. Si se filtra el token, no tiene valor por sí solo: habría que ir
a buscar los datos reales a la base ultra-protegida.

%======================================================================
\subsection{Estados del dato y encripción}
%======================================================================

Un mismo dato puede estar en distintos \textbf{estados}; las regulaciones definen
lineamientos para cada uno:
\begin{itemize}[nosep]
  \item \textbf{En reposo (at-rest):} el dato está quieto, no se procesa ni se mueve
        (base de datos, backup, Excel, bucket S3).
  \item \textbf{En tránsito / movimiento (in-motion / in-transit):} se mueve de una
        ubicación o estado a otro (a través de la red, un pendrive, una query).
  \item \textbf{En uso (in-use):} está siendo procesado/usado por una aplicación o CPU.
  \item \textbf{En memoria:} caso particular; puede verse como \emph{at-rest} en memoria
        esperando a la CPU, o \emph{in-transit} entre el bus de I/O y la CPU.
\end{itemize}

\subsubsection{Enforcement (Compiler-time)}
Asegura que el código cumpla políticas de seguridad \textbf{antes} de ejecutarse, vía
\textbf{análisis estático de código} en la compilación (ej.\ detectar que se imprimen
contraseñas/secretos). El compilador puede asegurar que datos confidenciales no se
filtren a zonas de memoria de menor seguridad. Ejemplo: \texttt{GuardedString} de Java
(guarda el string \textbf{encriptado} en memoria; un \texttt{String} clásico se ve en
texto plano inspeccionando el proceso).

\subsubsection{Encripción at-rest}
Objetivo: proteger datos inactivos (disco, pendrive, etc.). Mecanismos:
\begin{itemize}[nosep]
  \item \textbf{HSM} (\emph{Hardware Security Module}, FIPS 1--4): hardware
        \textbf{anti-tampering} que guarda secretos (ej.\ claves simétricas) y ofrece
        las operaciones de cifrado/descifrado \textbf{sin exponer la clave}.
  \item \textbf{Keystore}: software que guarda todas las claves encriptadas con una
        clave \textbf{máster}.
  \item \textbf{Esquema KEK-DEK (encripción por capas):}
        \begin{itemize}[nosep]
          \item \textbf{DEK} = \emph{Data Encryption Key} (cifra los datos).
          \item \textbf{KEK} = \emph{Key Encryption Key} (cifra la DEK).
        \end{itemize}
        Más seguro y flexible: se puede \textbf{rotar la KEK} sin re-encriptar los
        datos (la DEK no cambia). Más difícil de administrar (segunda capa).
        Ejemplos: EncFS/Loop-AES, EFS (Windows), FileVault (Mac OS X). Problemas
        asociados: \emph{Data Remanence} y \emph{Secure Deletion}.
\end{itemize}

\paragraph{At-rest en la nube.} La nube es infraestructura de terceros. Espectro de
opciones (de mayor control/menor costo a mayor eficiencia operativa/mayor costo):
\begin{center}
\begin{tabular}{@{}>{\bfseries}l p{10.5cm}@{}}
\toprule
Opción & Quién controla las claves \\
\midrule
Own Encryption & Encripción del lado del proceso \textbf{antes} de almacenar.
  \textbf{Independiente} del proveedor. \\
Hold Your Own Keys (HYOK) & El cliente gestiona las claves, posee las privadas y decide
  la rotación. \\
Bring Your Own Keys (BYOK) & El cliente trae sus claves. \\
Customer Managed Keys & El cliente decide rotación y puede encriptar lo que desee, pero
  \textbf{no accede} a las claves privadas. \\
Service Managed Keys & El \textbf{proveedor} decide la rotación; el cliente no accede a
  las claves ni las usa para otros contenidos. Máxima eficiencia, mínimo control. \\
\bottomrule
\end{tabular}
\end{center}
\textbf{Ojo con la ubicación de los datos:} distintas regiones/países aplican distintas
regulaciones.

\subsubsection{Encripción in-transit}
\begin{itemize}[nosep]
  \item Encriptar \textbf{antes} de mover (ej.\ cifrar un archivo con GPG antes de
        mandarlo por mail).
  \item Usar \textbf{protocolos de transporte} que aseguren cifrado: TLS (HTTPS),
        SMTP sobre TLS, IPSEC, SSH. (Acuerdan clave por asimétrica y luego cifran todo
        el tráfico de forma simétrica.)
\end{itemize}

\subsubsection{Enforcement (Execution-time)}
Técnica más usada para monitorear/controlar cómo se transmite la información entre
componentes. Se implementa a nivel \textbf{aplicativo} (buenos principios de seguridad,
no exponer secretos) o a nivel de \textbf{infraestructura} con \textbf{DLP} (\emph{Data
Loss Prevention}). Requiere entender el movimiento de datos (entre procesos, niveles de
seguridad o estados).

\subsubsection{Encripción en memoria (TME / TME-MK)}
En ambientes con múltiples VMs accediendo a la misma memoria hay riesgo de acceder a
información confidencial. Intel (y otros) implementaron en hardware:
\begin{itemize}[nosep]
  \item \textbf{TME} (\emph{Total Memory Encryption}): \textbf{una sola clave} compartida
        por todos los procesos. La desencriptación se hace \textbf{dentro de la CPU,
        antes del caché} (AES-XTS en los controladores DRAM/NVRAM).
  \item \textbf{TME-MK} (\emph{Multi-Key}): \textbf{múltiples claves}, con control de
        acceso soportado por los registros de \textbf{VT-x} (cada VM usa un \emph{KeyID}).
\end{itemize}

\subsubsection{SoC / IoT / Embedded Security}
Escenario de los más desfavorables: chips expuestos físicamente, arquitectura abierta,
\textbf{debug} accesible, conectividad \textbf{wireless}, gran \textbf{escala} y
\textbf{riesgo físico} (tampering). Se recomienda incluir un HSM embebido (ej.\ chip
ATECC608A). Capas: Perception, Application, Network.

%======================================================================
\subsection{Cumplimiento, detección y respuesta}
%======================================================================

\subsubsection{Comply (Cumplimiento)}
\prof{No basta con ser bueno, también hay que parecerlo.}
Hay que \textbf{demostrar} (ej.\ ante una auditoría) que la protección efectivamente
ocurre $\Rightarrow$ todo se \textbf{loguea} y queda evidencia.

\textbf{Retención:} guardar información que ya no es necesaria, por temas regulatorios /
legales / forenses. Ejemplos: HIPAA \textbf{6 años} desde el último uso; SOX \textbf{7
años}. Algunas leyes obligan a \textbf{rotar las claves de encripción} (ej.\ cada 10
años) y re-encriptar todo.

\subsubsection{Detect -- DLP (Data Loss Prevention)}
Herramientas que ayudan a clasificar, inventariar y detectar datos en todos sus estados.
\textbf{No son preventivas ni correctivas, sino DETECTIVAS}: se activan al fallar la
aplicación de una política (generan alerta). Son \textbf{complejos de implementar}; hoy
sus funcionalidades vienen incorporadas en productos de uso diario (pueden detectar,
ej.\ encripción no autorizada típica de \emph{ransomware}).

\subsubsection{Respond (Respuesta a incidentes)}
Equipos dedicados con protocolo. \textbf{Primeras 24 h:} registrar fecha/hora, alertar,
asegurar el lugar, \textbf{frenar la exfiltración}, documentar todo, entrevistar a los
involucrados, traer equipo \textbf{forense}, notificar a las autoridades. \textbf{Más
allá de 24 h:} trabajar con forenses, identificar obligaciones legales, reportar a la
dirección, identificar problemas futuros y \textbf{arreglar el problema}.
\begin{destacado}
\textbf{No ocultar} el incidente, pero comunicar \textbf{solo lo mínimo y necesario}.
Si hay regulación de datos personales de por medio, hay que informar al organismo (y,
según el caso, a bancos y clientes afectados).
\end{destacado}

%======================================================================
\subsection{Encripción Homomórfica}
%======================================================================

\begin{definicion}{Encripción homomórfica}
Esquema de cifrado que permite \textbf{operar sobre los datos cifrados sin
descifrarlos}. Formalmente, para una función $f$:
\[
f\big(\mathrm{Enc}_k(x)\big) = \mathrm{Enc}_k\big(f(x)\big).
\]
Es decir, operar sobre el texto cifrado y luego descifrar da el mismo resultado que
operar sobre el texto plano.
\end{definicion}

\paragraph{Utilidad.} Permite \textbf{computación en la nube preservando la privacidad}:
un tercero (proveedor) puede procesar/analizar datos cifrados \textbf{sin verlos nunca}
en claro. Útil cuando se quiere delegar cómputo sobre datos sensibles sin exponerlos.

\begin{ojo}
La diapositiva del teórico escribe la propiedad como $f(x)=y \Rightarrow f(\mathrm{Enc}(x))=f(\mathrm{Dec}(y))$,
pero la formulación conceptual correcta y la que conviene recordar es
$f(\mathrm{Enc}_k(x)) = \mathrm{Enc}_k(f(x))$: \textbf{operar sobre lo cifrado equivale a
cifrar el resultado de operar sobre lo plano}.
\end{ojo}

\end{document}
